% ============================================================================
% Fichier  : partie1/P01_C02_risque_cyber_incertitude.tex
% Titre    : Le Risque Cyber : Notions, Mesure, Incertitude
% Auteur   : MINKA MI NGUIDJOÏ Thierry Emmanuel
% Version  : 1.0 -- novembre 2026
% Corrigés : corriges/Corr_C02.tex
% ============================================================================
% Faits du cas (bible des cas) : BANQUE SANAGA (C01, P00) ; R-AG-01 ; FAIR
% 0,2 à 0,6 événement par an ; coût de l'incident 410 M FCFA, 3 jours ;
% prestataire DIGICORE sans clause d'audit ; aucun contrôle sur les
% changements (GCI). LiZbiZ : registre R1 à R8 et appétence (ACSI C03,
% repris sans modification), constats de l'audit des quatre centres.
% Faits nouveaux (reportés dans la bible) : registre 2025 de la SANAGA, six
% scénarios R-DC-01 à R-DC-06 ; revue des 12 cotes : 6 OUI / 2 NON / 4 NEK ;
% appétence provisoire adoptée en novembre 2026 (convention du corpus).
% Échelles : convention du corpus (ACSI, Tableau 3.2), précisée ici par des
% bornes de fréquence. Convention de réduction par étape du scénario :
% propre à SEC08.
% Sources : ISO 31000:2018 (définition reformulée) ; ISO/IEC 27005:2022 ;
% guide EBIOS RM (ANSSI, 2018) ; Cox (2008) sur les matrices de risque.
% ============================================================================

\chapter{Le Risque Cyber : Notions, Mesure, Incertitude}
\label{chap:risque}
\chaptermark{Le risque cyber}

\epigraph{%
    « Une cote sans preuve est une opinion\\
    qui a pris l'apparence d'un chiffre. »%
}{Principe du cours}

\niveauChapitre{Fondamental}{%
    Maîtrise du Chapitre~\ref{chap:grc}%
}

\begin{boxobjectifs}
\begin{itemize}
    \item Définir risque, actif, menace, vulnérabilité, source de risque,
          objectif visé, événement redouté et scénario.
    \item Expliquer ce qui distingue le risque cyber des autres risques
          opérationnels.
    \item Coter vraisemblance, impact et criticité sur une échelle commune,
          et en connaître les limites.
    \item Distinguer risque inhérent, risque résiduel et appétence, et ne
          réduire un risque que par des mesures testées.
    \item Distinguer l'incertitude aléatoire de l'incertitude épistémique.
    \item Reconnaître une cote sans preuve comme une incertitude (\NEK{}) et
          mesurer la base de preuve d'un registre par $\mu$, $U$ et $H$.
\end{itemize}
\end{boxobjectifs}

\bigskip

% ============================================================================
\section*{Avant de commencer : sur quoi repose ce « 1 »~?}
% ============================================================================

Douala, novembre 2026. Mme~ABENA relit le registre des risques que le
comité des risques avait approuvé un an plus tôt\footnote{Comme dans tout
le cours, la BANQUE SANAGA, ses dirigeants, les dates et les montants sont
fictifs~; les normes et les méthodes citées sont exactes.}. Il compte six
scénarios, tous sur le centre de données. L'un d'eux, « intrusion par
Internet », porte une vraisemblance de 1 sur 4~: rare. La RSSI demande au
responsable de l'exploitation d'où vient ce chiffre.

« Nous n'avons jamais été attaqués », répond-il. Elle lui demande alors qui
lit les journaux du pare-feu d'accès à Internet. Personne~: installé au
siège, hors du périmètre certifié, il conserve ses journaux trente jours,
puis les efface~; la journalisation du centre de données ne voit que ce
qui a franchi ce pare-feu. La banque ne savait pas si elle avait été
attaquée. Le « 1 » ne mesurait pas une fréquence d'attaques~; il mesurait
l'absence de regard.

\medskip

Le prologue a montré qu'un risque jamais apprécié n'est pas faible, mais
inconnu. Ce chapitre va plus loin~: même un risque apprécié et coté peut
reposer sur une inconnue. Il définit le vocabulaire du risque cyber, montre
comment on le cote, puis apprend à coter la cote elle-même.

\begin{boxalert}
\textbf{Violation typique.} Coter la vraisemblance d'une attaque à partir
de l'absence d'incident connu, alors que rien n'aurait permis de détecter
l'incident.

\smallskip
\textbf{Règle~:} l'absence de preuve d'attaque n'est une preuve d'absence
que si l'on regardait. Une cote fondée sur un silence non surveillé est
\NEK{}.
\end{boxalert}

\separateur

% ============================================================================
\section{Définir le Risque Cyber}
\label{sec:c02-definir}
% ============================================================================

\subsection{L'effet de l'incertitude sur les objectifs}

\begin{boxdef}
\termecle{Risque}\index{risque}~: effet de l'incertitude sur les objectifs.
L'effet est un écart, positif ou négatif, par rapport à ce qui est
attendu. On décrit un risque par ses sources, les événements possibles,
leurs conséquences et leur vraisemblance~\cite{iso31000_2018}.
\end{boxdef}
{\small\itshape Définition reformulée d'après l'ISO~31000:2018.}

\medskip

Un risque se rattache toujours à un \emph{objectif}. « Les postes de
Ngaoundéré sont chiffrés » n'est un risque qu'au regard d'un objectif, par
exemple «~servir les clients au guichet tous les jours ouvrés~». C'est
pourquoi le conseil, qui fixe les objectifs, fixe aussi le niveau de risque
acceptable (Chapitre~\ref{chap:grc}).

\begin{boxdef}
\termecle{Risque cyber}\index{risque cyber}~: risque dont la source est
une action, malveillante ou accidentelle, sur les systèmes d'information
ou les données, et dont les conséquences touchent les activités de
l'organisation~: interruption, perte financière, fuite d'information,
sanction, atteinte à la réputation.
\end{boxdef}

\subsection{Les deux portes d'entrée de l'ISO/IEC~27005}

L'ISO/IEC~27005:2022, guide de la gestion des risques de sécurité de
l'information, admet deux manières d'identifier les
risques~\cite{iso27005_2022}~: à partir des \emph{actifs} (on recense ce
qui a de la valeur, puis les menaces et les vulnérabilités qui le
concernent), ou à partir des \emph{événements} (on part des scénarios
stratégiques, des sources de risque et de leurs conséquences sur
l'activité). La première est exhaustive mais lourde~; la seconde va
droit aux scénarios qui comptent, au risque d'en oublier. Le
Chapitre~\ref{chap:iso27005} les met en œuvre~; ce chapitre en pose le
vocabulaire commun.

% ============================================================================
\section{Les Composantes d'un Risque}
\label{sec:c02-composantes}
% ============================================================================

\subsection{Huit notions}

\begin{boxdef}
\termecle{Valeur métier}\index{valeur métier} (ou actif essentiel)~:
activité, processus ou information qui a de la valeur pour l'organisation
(le service au guichet, le fichier des clients).

\medskip
\termecle{Bien support}\index{bien support} (ou actif support)~:
composant sur lequel repose une valeur métier~: matériel, logiciel,
réseau, local, personne, prestataire.

\medskip
\termecle{Menace}\index{menace}~: cause potentielle d'un incident (un
rançongiciel, une erreur d'administration, une coupure électrique).

\medskip
\termecle{Vulnérabilité}\index{vulnérabilité}~: faiblesse d'un bien
support qu'une menace peut exploiter (un partage de fichiers ouvert à
tous, un compte sans second facteur).

\medskip
\termecle{Source de risque}\index{source de risque}~: élément, personne ou
groupe susceptible d'engendrer le risque~; pour une attaque, l'attaquant
et ses capacités.

\medskip
\termecle{Objectif visé}\index{objectif visé}~: ce que la source de
risque cherche à obtenir (un gain financier, de l'information, une
nuisance).

\medskip
\termecle{Événement redouté}\index{événement redouté}~: atteinte à une
valeur métier, exprimée du point de vue de l'organisation (indisponibilité
des guichets pendant plusieurs jours).

\medskip
\termecle{Scénario}\index{scénario de risque}~: enchaînement qui mène
une source de risque, par des biens supports et leurs vulnérabilités,
jusqu'à l'événement redouté.
\end{boxdef}

Les notions de valeur métier, bien support, source de risque, objectif
visé et événement redouté sont celles de la méthode EBIOS Risk Manager,
publiée par l'Agence nationale de la sécurité des systèmes d'information
(\sigle{ANSSI}, France)~\cite{anssi2018ebios}. Le cours les retient parce
qu'elles séparent nettement ce que l'on protège (la valeur métier) de ce
par quoi on est attaqué (les biens supports). D'autres textes emploient
d'autres mots pour les mêmes idées~; le Tableau~2.1 donne les plus
courants.

\begin{tableaucours}{L{3.4cm}L{4.0cm}L{4.8cm}}{%
    \tableauHeader{Ce cours} &
    \tableauHeader{EBIOS Risk Manager} &
    \tableauHeader{Autres formulations courantes}
}
Valeur métier & Valeur métier & Actif essentiel, actif primaire \\
Bien support & Bien support & Actif support \\
Source de risque & Source de risque (SR) & Attaquant, menace au sens
large \\
Objectif visé & Objectif visé (OV) & Motivation de l'attaquant \\
Événement redouté & Événement redouté & Événement, conséquence \\
Scénario & Scénario stratégique, puis opérationnel & Scénario de risque,
scénario d'incident \\
Impact & Gravité & Conséquence \\
\end{tableaucours}
\vspace{-0.8em}
\begin{center}{\small\itshape Tableau~2.1. Correspondance des
vocabulaires (simplifiée)}\end{center}

\subsection{Le scénario R-AG-01, décomposé}

\begin{figure}[H]
\centering
\begin{tikzpicture}[
    etape/.style={draw=CouleurPrimaire, thick, rounded corners,
                  fill=CouleurPrimaire!8, text width=2.6cm,
                  minimum height=1.9cm, align=center, font=\scriptsize},
    fleche/.style={->, thick, CouleurSecondaire}]
    \node[etape] (S) at (0,0) {\textbf{Source de risque}\\groupe
        cybercriminel\\[2pt]\textbf{Objectif visé}\\extorsion};
    \node[etape] (B) at (3.3,0) {\textbf{Biens supports}\\messagerie,
        postes, partage de fichiers, comptes};
    \node[etape] (E) at (6.6,0) {\textbf{Événement redouté}\\guichets
        indisponibles plusieurs jours};
    \node[etape] (V) at (9.9,0) {\textbf{Valeur métier}\\service au
        guichet des agences};
    \draw[fleche] (S) -- (B);
    \draw[fleche] (B) -- (E);
    \draw[fleche] (E) -- (V);
    \node[font=\scriptsize, align=center, text width=6cm] at (3.3,-1.6)
        {vulnérabilités~: pas de second facteur,\\partage ouvert à toutes
        les agences, sauvegardes non testées};
\end{tikzpicture}
\caption{Les composantes du scénario R-AG-01}
\label{fig:c02-scenario}
\end{figure}

La décomposition dit où agir. On peut viser la source (on ne la contrôle
pas), les vulnérabilités des biens supports (second facteur,
cloisonnement), ou les conséquences sur la valeur métier (sauvegardes
restaurables, procédure de guichet en mode dégradé). Les mesures se
classent ainsi par \emph{étape du scénario}~: entrée, propagation,
récupération. Cette classification servira à coter le risque résiduel
(Section~\ref{sec:c02-residuel}).

% ============================================================================
\section{Ce qui rend le Risque Cyber Particulier}
\label{sec:c02-specificites}
% ============================================================================

Le risque cyber se gère avec les outils des risques opérationnels, mais il
en déjoue plusieurs hypothèses. Chacune a une conséquence sur la manière
de le coter.

\begin{tableaucours}{L{3.0cm}L{4.8cm}L{4.4cm}}{%
    \tableauHeader{Particularité} &
    \tableauHeader{Ce qu'elle signifie} &
    \tableauHeader{Conséquence pour la cote}
}
Un adversaire qui s'adapte & L'attaquant choisit sa cible et contourne
les mesures connues~; la fréquence dépend de la défense & Une fréquence
passée ne prédit pas la fréquence future \\
Peu d'historique & Les incidents sont rares, mal déclarés, et changent de
forme chaque année & Les cotes reposent souvent sur l'avis d'expert \\
Des incidents invisibles & Un attaquant peut rester des semaines dans un
réseau sans être détecté & L'absence d'incident connu ne vaut que si l'on
surveillait \\
Des effets simultanés & Un même scénario touche d'un coup de nombreux
sites (9 agences en une nuit) & L'impact se cote sur le scénario entier,
pas par site \\
Des pertes très dispersées & La plupart des incidents coûtent peu~;
quelques-uns coûtent énormément & Une moyenne trompe~; il faut un
intervalle (Chapitre~\ref{chap:fair}) \\
\end{tableaucours}
\vspace{-0.8em}
\begin{center}{\small\itshape Tableau~2.2. Cinq particularités du risque
cyber}\end{center}

\begin{boxinfo}
\textbf{Un risque qui dépend de la défense.} Une coupure électrique ne
devient pas plus fréquente parce que la banque a installé un groupe
électrogène. Une attaque, si~: quand une cible se protège, l'attaquant en
cherche une autre, ou une autre porte. La vraisemblance d'un scénario cyber
n'est donc jamais une donnée de la nature~; c'est une hypothèse sur un
adversaire, qui se révise.
\end{boxinfo}

% ============================================================================
\section{Coter un Risque}
\label{sec:c02-coter}
% ============================================================================

\subsection{Une échelle commune}

Pour comparer des risques de natures différentes, et pour que le conseil
puisse arbitrer, on les cote sur une même échelle. Le cours reprend la
convention à quatre niveaux du corpus, et la précise par des bornes de
fréquence~; chaque organisation fixe la sienne.

\begin{tableaucours}{cL{4.2cm}L{7.4cm}}{%
    \tableauHeader{Niv.} &
    \tableauHeader{Vraisemblance} &
    \tableauHeader{Impact (le critère le plus défavorable l'emporte)}
}
1 & Rare~: moins d'une fois en 10~ans & Mineur~: perte inférieure à
5~millions FCFA, interruption de moins de 4~heures \\
2 & Possible~: une fois en 3 à 10~ans & Modéré~: de 5 à 50~millions FCFA,
interruption de moins de 24~heures \\
3 & Probable~: une fois en 1 à 3~ans & Majeur~: de 50 à 500~millions FCFA,
interruption de 1 à 3~jours, réserve du commissaire aux comptes ou compte
rendu exigé par le régulateur \\
4 & Quasi certain~: au moins une fois par an & Critique~: plus de
500~millions FCFA, interruption de plus de 3~jours, sanction d'un
régulateur \\
\end{tableaucours}
\vspace{-0.8em}
\begin{center}{\small\itshape Tableau~2.3. Échelles de cotation
(convention du cours)}\end{center}

\begin{boxdef}
\termecle{Criticité}\index{criticité} (ou score)~: produit
$\text{vraisemblance} \times \text{impact}$, de 1 à 16. Elle sert à
classer les risques et à les situer par rapport à l'appétence.
\end{boxdef}

L'incident de septembre 2026 a coûté 410~millions FCFA et fermé les
guichets trois jours~: il se situe au niveau~3 d'impact. Le scénario
redouté de R-AG-01 va plus loin~: une propagation à l'ensemble du réseau,
au-delà de trois jours. Son impact se cote~4. On cote le scénario que l'on
redoute, pas seulement celui que l'on a vécu.

\subsection{L'appétence, traduite sur l'échelle}

En novembre 2026, le comité des risques du conseil adopte, à titre
provisoire, la règle d'appétence commune au corpus~:
\begin{itemize}
    \item une criticité résiduelle de \textbf{9 ou plus} est
          \emph{intolérable} et doit être traitée sans délai~;
    \item de \textbf{5 à 8}, le risque est \emph{traité dans l'année}~;
    \item de \textbf{1 à 4}, il est \emph{accepté sous surveillance}, par
          une décision écrite du propriétaire~;
    \item tout risque d'impact~4 est présenté au comité, quelle que soit sa
          criticité.
\end{itemize}

La dernière règle n'est pas un détail. Elle corrige le principal défaut
des matrices, qu'il faut maintenant examiner.

\subsection{Ce que la matrice cache}

La matrice vraisemblance~$\times$ impact est l'outil le plus répandu de la
gestion des risques. Elle est aussi l'un des plus critiqués. Cox a montré
qu'elle peut classer deux risques dans le mauvais ordre, et donner la même
couleur à des risques dont les pertes attendues diffèrent de plusieurs
ordres de grandeur~\cite{cox2008matrices}. Trois raisons~:
\begin{itemize}
    \item \textbf{des échelles ordinales}~: le niveau~4 n'est pas « deux
          fois » le niveau~2~; leurs bornes (une fois par an, une fois en
          trois ans) sont dans un rapport de 1 à 3 au moins, et les bornes
          d'impact dans un rapport de 1 à 10~;
    \item \textbf{un produit qui écrase}~: un risque fréquent et mineur
          (4~$\times$~1) et un risque rare et critique (1~$\times$~4) ont la
          même criticité, 4~;
    \item \textbf{des frontières arbitraires}~: une fréquence de 0,32 par
          an et une de 0,34 tombent dans deux niveaux différents, une
          fréquence de 0,34 et une de 0,99 dans le même.
\end{itemize}

\begin{boxalert}
\textbf{Violation typique.} Accepter un risque rare et critique parce que
sa criticité tombe en zone verte.

\smallskip
\textbf{Règle~:} la criticité classe, elle ne décide pas. Un risque
d'impact critique se présente toujours à l'instance qui répond des
risques, avec son scénario. Quand la décision dépend de l'ordre de
grandeur des pertes, on quantifie (Chapitre~\ref{chap:fair}).
\end{boxalert}

% ============================================================================
\section{Inhérent, Résiduel, Appétence}
\label{sec:c02-residuel}
% ============================================================================

\subsection{Trois niveaux}

\begin{boxdef}
\termecle{Risque inhérent}\index{risque inhérent}~: niveau du risque
avant la prise en compte des mesures de sécurité.

\medskip
\termecle{Risque résiduel}\index{risque résiduel}~: niveau du risque qui
subsiste compte tenu des mesures réellement en place \emph{et dont
l'efficacité a été établie}.

\medskip
\termecle{Appétence au risque}\index{appétence au risque}~: niveau de
risque que l'organisation accepte pour atteindre ses objectifs, fixé par
l'organe de gouvernance (Chapitre~\ref{chap:grc}).
\end{boxdef}

\subsection{Une convention de réduction}

Pour que deux analystes obtiennent le même résiduel, il faut une règle de
réduction écrite. Le cours retient la suivante, fondée sur les étapes du
scénario.

\begin{boxprocedure}
\textbf{Réduire un risque, étape par étape.}
\begin{itemize}
    \item On classe les mesures selon l'étape du scénario sur laquelle
          elles agissent~: \emph{entrée} (empêcher l'accès initial),
          \emph{propagation} (empêcher l'extension), \emph{récupération}
          (limiter les conséquences).
    \item Une étape couverte par au moins une mesure \textbf{testée et
          efficace} abaisse d'un niveau~: la vraisemblance pour l'entrée et
          la propagation, l'impact pour la récupération.
    \item Une étape ne compte qu'une fois, quel que soit le nombre de
          mesures qui la couvrent~; aucun niveau ne descend sous~1.
    \item Une mesure non testée, ou cotée \NEK{}, n'abaisse rien.
    \item Le papier de travail cite, pour chaque réduction, la mesure et la
          preuve de son test.
\end{itemize}
\end{boxprocedure}

\subsection{R-AG-01, coté}

\begin{boxscenario}{Opération SANAGA. R-AG-01, inhérent et résiduel}
\textbf{Inhérent.} Une banque de 38 agences qui reçoit chaque jour des
courriels de l'extérieur sera visée au moins une fois par an~:
vraisemblance~4. Le scénario redouté touche tout le réseau plus de trois
jours~: impact~4. Criticité inhérente~: $4 \times 4 = 16$.

\medskip
\textbf{Mesures de 2025}, avec les cotes de l'audit interne (prologue,
Tableau~P.3)~:
\begin{itemize}
    \item entrée~: filtrage des pièces jointes et antivirus, \OUI{}~; second
          facteur, \NON{} $\Rightarrow$ étape couverte, $-1$ en
          vraisemblance~;
    \item propagation~: cloisonnement, \NON{} $\Rightarrow$ aucune
          réduction~;
    \item récupération~: sauvegardes hors ligne testées, \NON{}~; plan de
          réponse testé, \NON{} $\Rightarrow$ aucune réduction.
\end{itemize}

\textbf{Résiduel.} Vraisemblance~3, impact~4~: criticité $3 \times 4 = 12$,
intolérable.

\medskip
\textbf{Analyse} (confiance moyenne). Le registre de 2025 aurait dû porter
un risque intolérable sur les agences. Il n'en portait aucun, parce que le
scénario n'y figurait pas. La cote résiduelle de 12 n'est pas une
prédiction~: elle dit que, faute de mesures sur deux étapes sur trois, la
banque dépendait d'une seule barrière, et que cette barrière arrête les
attaques connues, pas les nouvelles.
\end{boxscenario}

% ============================================================================
\section{Une Cote sans Preuve est une Incertitude}
\label{sec:c02-incertitude}
% ============================================================================

\subsection{Deux incertitudes}

Le mot « incertitude » figure dans la définition du risque. Il recouvre
deux choses que la gestion des risques confond souvent.

\begin{boxdef}
\termecle{Incertitude aléatoire}\index{incertitude aléatoire}~: part de
l'incertitude qui tient au hasard des événements~: même parfaitement
informé, on ne sait pas si l'attaque aura lieu cette année. Elle ne se
réduit pas~; elle se décrit par une probabilité.

\medskip
\termecle{Incertitude épistémique}\index{incertitude épistémique}~: part
de l'incertitude qui tient à ce que l'on ne sait pas de l'état du système
ou de la menace~: on ignore si les journaux contiennent des traces
d'intrusion. Elle se réduit en collectant des preuves.
\end{boxdef}

\begin{tableaucours}{L{2.8cm}L{4.6cm}L{4.8cm}}{%
    \tableauHeader{} &
    \tableauHeader{Aléatoire} &
    \tableauHeader{Épistémique}
}
Porte sur & Les événements futurs & L'état présent du système et de la
menace \\
Se réduit~? & Non & Oui, par des preuves \\
Se mesure par & Une probabilité, un intervalle de pertes & Le taux
d'incertitude $U$ et l'entropie $H$ des réponses \\
À la BANQUE SANAGA & Le groupe attaquera-t-il en 2027~? & Les journaux du
pare-feu contiennent-ils des traces d'intrusion~? \\
\end{tableaucours}
\vspace{-0.8em}
\begin{center}{\small\itshape Tableau~2.4. Incertitude aléatoire et
incertitude épistémique}\end{center}

Une vraisemblance cotée « 1 » sans base de preuve mélange les deux~: elle
prétend décrire le hasard (l'aléatoire), alors qu'elle traduit une
ignorance (l'épistémique). Le programme de l'auditologie fait de cette
distinction l'une de ses thèses~\cite{minka2026auditologie,minka2026nek}~:
avant de mesurer un risque, on mesure ce que l'on sait de lui.

\subsection{Coter la cote}

Chaque cote d'un registre repose sur une \termecle{base de
preuve}\index{base de preuve}~: des mesures internes (journaux, incidents
recensés, tests), des données externes (statistiques du secteur, alertes
des autorités), un avis d'expert motivé, ou rien. On cote donc la cote
elle-même~:
\begin{itemize}
    \item \OUI{}~: la base de preuve existe et soutient la cote (un avis
          d'expert motivé et écrit compte, et se signale comme tel)~;
    \item \NON{}~: la base de preuve existe et contredit la cote~;
    \item \NEK{}~: aucune base de preuve, ou une base qui ne permet pas de
          conclure (une déclaration, un silence non surveillé).
\end{itemize}

\begin{boxscenario}{Opération SANAGA. Le registre de 2025, revu}
Mme~ABENA revoit les douze cotes (six vraisemblances, six impacts) des six
scénarios du registre approuvé en novembre 2025.

\medskip
\begin{tableaucours}{cL{3.4cm}ccL{3.4cm}cc}{%
    \tableauHeader{Réf.} &
    \tableauHeader{Scénario} &
    \tableauHeader{V} &
    \tableauHeader{I} &
    \tableauHeader{Base de la vraisemblance} &
    \tableauHeader{V} &
    \tableauHeader{I}
}
DC-01 & Coupure électrique prolongée & 2 & 4 & Journal des coupures sur
cinq ans & \OUI{} & \OUI{} \\
DC-02 & Défaillance du prestataire DIGICORE & 2 & 3 & « Prestataire
fiable »~; aucune revue, pas de droit d'audit & \NEK{} & \OUI{} \\
DC-03 & Intrusion par Internet & 1 & 4 & « Jamais attaqués »~; journaux
du pare-feu d'accès non lus & \NEK{} & \OUI{} \\
DC-04 & Erreur lors d'un changement & 1 & 3 & Sept incidents liés à des
changements en 2025 & \NON{} & \OUI{} \\
DC-05 & Copie de données par un administrateur & 1 & 4 & « Confiance dans
les équipes » & \NEK{} & \NEK{} \\
DC-06 & Déni de service des services en ligne & 1 & 3 & Quatre attaques
signalées par le fournisseur d'accès en 2025 & \NON{} & \OUI{} \\
\end{tableaucours}

\smallskip
{\small Colonnes~: V et I cotés (1 à 4)~; puis cote de la base de preuve
de la vraisemblance (V) et de l'impact (I). Les impacts reposent sur
l'analyse d'impact sur l'activité de 2024, sauf pour DC-05~: personne n'a
établi quelles données un administrateur pouvait copier.}
\end{boxscenario}

\begin{boxlecture}{Base de preuve du registre de 2025}
\textbf{Constatation.} Sur douze cotes~: 6~\OUI{}, 2~\NON{}, 4~\NEK{}.
\[
\mu = 0{,}50~;\quad U = \tfrac{4}{12} \approx 0{,}33~;\quad
H = -\left(\tfrac{1}{2}\log_2\tfrac{1}{2} + \tfrac{1}{6}\log_2\tfrac{1}{6}
  + \tfrac{1}{3}\log_2\tfrac{1}{3}\right) \approx 0{,}50 + 0{,}43 + 0{,}53
  \approx 1{,}46~\text{bit}.
\]
Par colonne~: impacts 5/0/1 ($\mu \approx 0{,}83$, $U \approx 0{,}17$,
$H \approx 0{,}65$)~; vraisemblances 1/2/3 ($\mu \approx 0{,}17$,
$U = 0{,}50$, $H \approx 1{,}46$).

\medskip
\textbf{Analyse} (confiance élevée). Le registre est solide sur les
conséquences et fragile sur les fréquences~: une seule vraisemblance sur
six est fondée. Les deux \NON{} sont les plus graves~: DC-04 et DC-06
étaient contredits par des pièces que la banque détenait. DC-04 annonçait
la bascule ratée de mars 2026. Enfin, ce tableau ne mesure que ce que le
registre contient~: le scénario des agences, absent, échappe même à cette
revue. C'est le \NEK{} du prologue, au niveau du périmètre.

\medskip
\hfill\lecture{0{,}50}{0{,}33}{1{,}46}
\end{boxlecture}

\subsection{Quand la cote est un intervalle}

Une estimation honnête de la vraisemblance prend souvent la forme d'un
intervalle. L'atelier de quantification du cours estime la fréquence de
R-AG-01 entre 0,2 et 0,6 événement par an (Chapitre~\ref{chap:fair}).
Sur l'échelle du Tableau~2.3, cet intervalle chevauche deux niveaux~: de 0,2
à 0,33, niveau~2~; de 0,33 à 0,6, niveau~3. Le résiduel vaut donc 8 ou 12
selon le niveau retenu~: « à traiter dans l'année » ou « intolérable ».

\begin{boxinfo}
\textbf{Une décision qui dépend de l'incertitude.} Quand l'intervalle
d'une cote traverse un seuil d'appétence, la décision dépend de ce que
l'on ignore. Deux réponses sont légitimes~: retenir le niveau le plus
défavorable en le signalant, ou réduire l'incertitude avant de décider
(une revue des journaux, une simulation d'attaque, une quantification).
Choisir silencieusement le niveau favorable ne l'est pas.
\end{boxinfo}

\subsection{La fiche de risque, avec sa base de preuve}

\begin{boxobjet}{R-DC-03 · Intrusion par Internet (fiche révisée)}
\textbf{Valeur métier~:} les opérations du système central (comptes,
virements).

\textbf{Biens supports~:} pare-feu, serveurs exposés, comptes
d'administration à distance.

\textbf{Source de risque et objectif visé~:} groupe cybercriminel, fraude
ou extorsion.

\textbf{Événement redouté~:} prise de contrôle du système central.

\textbf{Vraisemblance~:} 1 en 2025. Base~: absence d'incident connu,
journaux du pare-feu d'accès effacés à trente jours sans lecture. Cote de
la base~: \NEK{}.

\textbf{Impact~:} 4. Base~: analyse d'impact sur l'activité de 2024. Cote
de la base~: \OUI{}.

\textbf{Pour réduire l'incertitude~:} centraliser et conserver les
journaux un an~; faire analyser trois mois de journaux par un tiers~;
commander un test d'intrusion externe (UE MSN-2S3-SEC03).

\textbf{Décision provisoire~:} faute de fréquence observée, la
vraisemblance est recalculée par la convention de réduction~: inhérent~4
(serveurs exposés à Internet), entrée couverte par des mesures testées au
centre (second facteur, correctifs~: prologue, Tableau~P.3), $-1$~;
propagation couverte (cloisonnement testé), $-1$. Vraisemblance~2,
impact~4~: criticité~8, à traiter dans l'année, et présentée au comité en
raison de l'impact~4. À confirmer par la revue des journaux.
\end{boxobjet}

% ============================================================================
\section{Le Cadre de Référence}
\label{sec:c02-cadre}
% ============================================================================

L'article~26 de la loi n°~2010/012, cité au Chapitre~\ref{chap:grc}, impose
aux exploitants de systèmes d'information des dispositifs normalisés pour
\emph{identifier} et \emph{évaluer} en permanence les risques de
sécurité~\cite{loi2010012}. Pour une banque, le règlement COBAC R-2016/04
range les risques liés aux systèmes d'information parmi les risques que le
dispositif de contrôle interne doit couvrir~\cite{cobac2016r04}.

\begin{boxinfo}
\textbf{Analyse} (lecture du cours). Ni la loi ni le règlement ne disent
comment coter. Mais une évaluation que l'on ne peut pas justifier n'en est
pas une~: devant le superviseur comme devant un juge, une cote se défend
par sa base de preuve. Le registre qui distingue ce qu'il sait de ce qu'il
suppose est plus défendable que celui qui affiche des chiffres sans
source. Le Chapitre~\ref{chap:droit} revient sur ces obligations.
\end{boxinfo}

% ============================================================================
\section{Ce que la Cote demande à chaque Métier}
\label{sec:c02-metiers}
% ============================================================================

\begin{boxmetier}{Gestionnaire des risques cyber}
\textbf{Documenter chaque cote}
\begin{enumerate}
    \item Écrire, pour chaque vraisemblance et chaque impact, la base de
          preuve et sa date.
    \item Coter la base~: \OUI{}, \NON{} ou \NEK{}~; ne jamais transformer
          un \NEK{} en « faible ».
    \item Donner un intervalle quand l'estimation le justifie, et signaler
          tout intervalle qui traverse un seuil d'appétence.
    \item Ne réduire le résiduel que par des mesures testées, selon la
          convention écrite.
\end{enumerate}
\end{boxmetier}

\begin{boxmetier}{Auditeur en cybersécurité}
\textbf{Auditer le registre, pas seulement les mesures}

L'auditeur prélève des cotes et demande leur base. Une vraisemblance sans
base est une constatation~; une vraisemblance contredite par une pièce
détenue par l'organisation est une constatation grave, parce qu'elle
montre que l'information existait et n'a pas été utilisée.
\end{boxmetier}

\begin{boxmetier}{RSSI}
\textbf{Transformer les \NEK{} en travaux}

Chaque \NEK{} du registre appelle une action qui produit une preuve~:
lecture de journaux, test de restauration, revue d'un prestataire, test
d'intrusion. Le RSSI les classe par l'importance de la décision qu'ils
débloquent~: d'abord les cotes dont l'intervalle traverse un seuil.
\end{boxmetier}

\begin{boxmetier}{Responsable GRC}
\textbf{Une échelle, une convention}

Le responsable GRC fait approuver par le conseil l'échelle, l'appétence et
la convention de réduction, puis veille à ce que tous les registres de
l'organisation les appliquent. Deux registres qui cotent différemment ne
peuvent pas être comparés, donc pas arbitrés.
\end{boxmetier}

\separateur

\begin{boxresume}
\textbf{Ce qu'il faut retenir de ce chapitre~:}
\begin{itemize}
    \item \textbf{Un risque} est l'effet de l'incertitude sur un objectif~;
          un scénario relie une source de risque, par des biens supports
          vulnérables, à un événement redouté sur une valeur métier.
    \item \textbf{Le risque cyber} dépend d'un adversaire qui s'adapte, a
          peu d'historique, peut rester invisible, frappe plusieurs sites à
          la fois et produit des pertes très dispersées.
    \item \textbf{La criticité classe, elle ne décide pas}~: les matrices
          écrasent les écarts~; tout impact critique remonte à l'instance
          qui répond des risques.
    \item \textbf{Le résiduel} ne se réduit que par des mesures testées,
          selon une convention écrite~: une étape du scénario couverte
          abaisse d'un niveau.
    \item \textbf{Deux incertitudes}~: l'aléatoire se décrit, l'épistémique
          se réduit par des preuves.
    \item \textbf{On cote la cote}~: \OUI{}, \NON{} ou \NEK{} selon sa base
          de preuve. Le registre 2025 de la SANAGA~: $\mu = 0{,}50$,
          $U \approx 0{,}33$, et une seule vraisemblance fondée sur six.
\end{itemize}
\end{boxresume}

\begin{boxinfo}
\textbf{Points de contrôle.} Vous avez compris ce chapitre si vous savez~:
\begin{itemize}
    \item décomposer un scénario en source, objectif visé, biens supports,
          vulnérabilités, événement redouté et valeur métier~;
    \item calculer une criticité inhérente et résiduelle avec la convention
          de réduction~;
    \item montrer, sur un exemple, ce qu'une matrice de risques cache~;
    \item coter la base de preuve d'une vraisemblance, et dire ce qu'il
          faudrait pour la rendre \OUI{}.
\end{itemize}
\end{boxinfo}

% ============================================================================
\section*{Exercices}
% ============================================================================

\begin{boxexo}[title={\ding{45}\quad Exercice~2.1. Le vocabulaire du scénario \niveaubase}]
Associez chacun des éléments suivants, tirés de l'incident de septembre
2026, à l'une des notions du chapitre (valeur métier, bien support,
menace, vulnérabilité, source de risque, objectif visé, événement redouté,
scénario).
\begin{enumerate}[(a)]
    \item Un groupe cybercriminel organisé.
    \item Obtenir le paiement d'une rançon.
    \item Le partage de fichiers ouvert à toutes les agences.
    \item Les postes de travail des agences.
    \item Le service au guichet.
    \item La fermeture des guichets pendant plusieurs jours.
    \item Un rançongiciel.
    \item Courriel piégé, vol d'identifiants, déplacement par le partage,
          chiffrement des postes.
\end{enumerate}
\end{boxexo}

\begin{boxexo}[title={\ding{45}\quad Exercice~2.2. Six mois après \niveaumoyen}]
En mars 2027, la situation de R-AG-01 est la suivante~: le second facteur
est déployé et testé sur tous les comptes~; le partage de fichiers est
cloisonné par agence, mais le cloisonnement n'a pas encore été testé~; les
sauvegardes hors ligne ont été restaurées avec succès lors d'un test.
\begin{enumerate}[(a)]
    \item Calculez la criticité résiduelle avec la convention du chapitre
          (inhérent~: $4 \times 4$). Quelle décision l'appétence impose-t-elle~?
    \item Que deviendrait le résiduel si le cloisonnement était testé avec
          succès~?
    \item Le second facteur n'a rien changé au calcul. Expliquez pourquoi,
          et dites si c'est un défaut de la convention ou une qualité.
\end{enumerate}
\end{boxexo}

\begin{boxexo}[title={\ding{45}\quad Exercice~2.3. Ce que cache la matrice \niveaumoyen}]
Le registre de la banque contient deux risques~:
\begin{itemize}
    \item A~: erreurs de saisie au guichet, vraisemblance~4, impact~1~;
    \item B~: fraude massive au virement par compromission du système de
          paiement, vraisemblance~1, impact~4.
\end{itemize}
\begin{enumerate}[(a)]
    \item Calculez leur criticité. Que conclurait un comité qui lit
          seulement la criticité~?
    \item Les erreurs A se produisent douze fois par an et coûtent
          1~million FCFA chacune~; la fraude B a une chance sur vingt de se
          produire chaque année et coûterait 2~milliards FCFA. Calculez la
          perte annuelle attendue de chacun.
    \item Expliquez l'écart entre (a) et (b) à l'aide des trois défauts des
          matrices vus dans le chapitre.
    \item Quelle règle de l'appétence de la banque évite la mauvaise
          décision~?
\end{enumerate}
\end{boxexo}

\begin{boxexo}[title={\ding{45}\quad Exercice~2.4. La base de preuve du registre de LiZbiZ \niveauexpert}]
Mme~NKENG reprend le registre R1 à R8 de LiZbiZ (cotes de probabilité
inchangées) et recherche la base de chaque probabilité~:
\begin{itemize}
    \item R1, ajustement frauduleux du stock~: écart de 186~millions FCFA
          constaté, ajustements identifiés~;
    \item R2, fournisseurs fictifs~: fraude TRANSLOG établie~;
    \item R3, compte d'un agent parti~: « on désactive les comptes au
          départ », déclaration du responsable informatique, sans
          extraction~;
    \item R4, arrêt prolongé du centre de Garoua~: journal des coupures sur
          trois ans, un arrêt de plus d'un jour~;
    \item R5, perte de données faute de sauvegarde restaurable~: « les
          sauvegardes tournent chaque nuit », aucun test de restauration~;
    \item R6, fuite de données des clients~: aucune information~;
    \item R7, erreurs de facturation~: trois erreurs par mois relevées par
          le contrôle de gestion~;
    \item R8, rançongiciel (probabilité~2)~: « nous n'avons jamais été
          touchés »~; l'audit des quatre centres a pourtant relevé des
          comptes par défaut, des services en clair et des flux non
          autorisés.
\end{itemize}
\begin{enumerate}[(a)]
    \item Cotez chaque base \OUI{}, \NON{} ou \NEK{}, en justifiant.
    \item Calculez $\mu$, $U$ et $H$.
    \item Proposez une nouvelle cote de probabilité pour R8, recalculez sa
          criticité résiduelle (impact~4, aucune mesure testée) et dites ce
          que l'appétence impose.
    \item Rédigez pour Mme~NKENG une note de dix lignes~: quelles cotes
          rétablir en priorité, par quelles preuves, et pourquoi dans cet
          ordre.
\end{enumerate}
\end{boxexo}

\bigskip

\begin{boxinfo}
\textbf{Pour aller plus loin.}
\begin{itemize}
    \item L'ISO~31000:2018~\cite{iso31000_2018} et l'ISO/IEC~27005:2022~\cite{iso27005_2022},
          pour le vocabulaire et le processus.
    \item Le guide EBIOS Risk Manager~\cite{anssi2018ebios}, pour les notions
          de valeur métier, bien support et source de risque.
    \item L'article de Cox sur les matrices de risque~\cite{cox2008matrices}.
    \item L'article sur l'incertitude épistémique~\cite{minka2026nek}, qui
          fonde la cote \NEK{}.
\end{itemize}
\end{boxinfo}

\begin{boxinfo}
\textbf{Transition.} La Partie~I a posé qui décide (Chapitre~\ref{chap:grc})
et ce que l'on décide (ce chapitre). La Partie~II passe aux méthodes~:
comment conduire une appréciation des risques complète, de l'établissement
du contexte à l'acceptation du résiduel. C'est l'objet du
Chapitre~\ref{chap:iso27005}.
\end{boxinfo}
